\section*{Chapitre \ref{ch:b1:quantum-numbers} --- Nombres quantiques et configurations électroniques}

\begin{solution}{exo:b1:quantum-numbers:1}
(a) Impossible~: $l \le n - 1 = 1$. (b) Possible (un électron 3p).
(c) Impossible~: $m_s = \pm\tfrac12$. (d) Possible (un électron 4f).
(e) Impossible~: pour $l = 0$, $m_l$ ne peut valoir que 0.
\end{solution}

\begin{solution}{exo:b1:quantum-numbers:2}
3d~: 5 orbitales, 10 électrons~; 4f~: 7 orbitales, 14 électrons~; 5p~: 3
orbitales, 6 électrons~; couche $n = 2$~: $n^2 = 4$ orbitales, $2n^2 = 8$
électrons.
\end{solution}

\begin{solution}{exo:b1:quantum-numbers:3}
\ce{O}~: $1\text{s}^2 2\text{s}^2 2\text{p}^4 = [\ce{He}]\,2\text{s}^2 2\text{p}^4$,
6 \omterm{def:b1:quantum-numbers:configuration}{électrons de valence}. \ce{P}~: $1\text{s}^2 2\text{s}^2 2\text{p}^6 3\text{s}^2
3\text{p}^3 = [\ce{Ne}]\,3\text{s}^2 3\text{p}^3$, 5. \ce{K}~: $[\ce{Ar}]\,4\text{s}^1$,
1. \ce{Ca}~: $[\ce{Ar}]\,4\text{s}^2$, 2. \ce{Br}~: $[\ce{Ar}]\,3\text{d}^{10}
4\text{s}^2 4\text{p}^5$, 7 (la \omterm{def:b1:quantum-numbers:orbital}{sous-couche} 3d pleine appartient au cœur).
% ledger: gs:K, gs:Ca
\end{solution}

\begin{solution}{exo:b1:quantum-numbers:4}
\ce{N}, 2p~: \omorbs{u,u,u}, 3 célibataires. \ce{S}, 3p~: \omorbs{ud,u,u},
2 célibataires. \ce{Cl}, 3p~: \omorbs{ud,ud,u}, 1 célibataire. Dans
chaque cas, 2s ou 3s est pleine, \omorbs{ud}.
\end{solution}

\begin{solution}{exo:b1:quantum-numbers:5}
\ce{Na+}, \ce{Mg^{2+}}, \ce{Al^{3+}}, \ce{F-}, \ce{O^{2-}}~: $1\text{s}^2
2\text{s}^2 2\text{p}^6$, isoélectroniques du néon. \ce{Cl-}, \ce{S^{2-}},
\ce{Ca^{2+}}~: $[\ce{Ne}]\,3\text{s}^2 3\text{p}^6$, isoélectroniques de
l'argon.
\end{solution}

\begin{solution}{exo:b1:quantum-numbers:6}
Les électrons 4s partent les premiers (\cref{prop:b1:quantum-numbers:ions}).
\ce{Ti^{2+}}~: $[\ce{Ar}]\,3\text{d}^2$, 2 célibataires~; \ce{Mn^{2+}}~:
$3\text{d}^5$, 5~; \ce{Co^{2+}}~: $3\text{d}^7$, 3 (\omorbs{ud,ud,u,u,u})~;
\ce{Ni^{2+}}~: $3\text{d}^8$, 2~; \ce{Zn^{2+}}~: $3\text{d}^{10}$, 0.
% ledger: gs:Ti2+, gs:Mn2+, gs:Co2+, gs:Zn2+
\end{solution}

\begin{solution}{exo:b1:quantum-numbers:7}
$\Delta E = 13.6\,(1/9 - 1/16) = 13.6 \times 7/144 = \qty{0.661}{eV}$~;
$\lambda = 1240/0.661 = \qty{1876}{nm} \approx \qty{1.88}{\micro m}$~:
infrarouge (première raie de la série de Paschen).
\end{solution}

\begin{solution}{exo:b1:quantum-numbers:8}
(a) $[\ce{Ne}]\,3\text{s}^2 3\text{p}^4$~: le soufre, $Z = 16$.
(b) $[\ce{Ar}]\,3\text{d}^{10} 4\text{s}^2 4\text{p}^5$~: le brome, $Z = 35$.
(c) $[\ce{Kr}]\,5\text{s}^1$~: le rubidium, $Z = 37$.
(d) $[\ce{Ar}]\,3\text{d}^3 4\text{s}^2$~: le vanadium, $Z = 23$.
Un électron 3p du soufre~: $(3, 1, -1, +\tfrac12)$, par exemple.
% ledger: gs:V
\end{solution}

\begin{solution}{exo:b1:quantum-numbers:9}
Pour 4s, $n + l = 4 + 0 = 4$~; pour 3d, $n + l = 3 + 2 = 5$~: 4s vient
en premier. Après 3p ($n + l = 4$), la \omterm{def:b1:quantum-numbers:orbital}{sous-couche} suivante est 4s
($n + l = 4$, $n$ plus grand), qui ouvre la période 4~; 3d ($n + l = 5$)
ne vient qu'après 4s. La période 3 ne contient donc que 3s et 3p~:
$2 + 6 = 8$ éléments.
\end{solution}

\begin{solution}{exo:b1:quantum-numbers:10}
Énergie d'ionisation $Z^2E_R = 4 \times 13.6 = \qty{54.4}{eV}$, quatre
fois celle de l'hydrogène. $2 \to 1$~: $\Delta E = 54.4\,(1 - 1/4) = \qty{40.8}{eV}$,
$\lambda = 1240/40.8 = \qty{30.4}{nm}$, quatre fois plus courte que la
raie de l'hydrogène à \qty{121.6}{nm}~: toutes les énergies sont
multipliées par $Z^2 = 4$.
\end{solution}

\begin{solution}{exo:b1:quantum-numbers:11}
(a) \omterm{def:b1:quantum-numbers:energy-level}{État excité} du béryllium (\omterm{def:b1:quantum-numbers:energy-level}{état fondamental} $1\text{s}^2 2\text{s}^2$).
(b) Impossible~: 2p contient au plus 6 électrons. (c) \omterm{def:b1:quantum-numbers:energy-level}{État excité} du
béryllium. (d) \omterm{def:b1:quantum-numbers:energy-level}{État fondamental} du phosphore. (e) Impossible~: une
\omterm{def:b1:quantum-numbers:orbital}{sous-couche} s contient deux électrons (Pauli).
\end{solution}

\begin{solution}{exo:b1:quantum-numbers:12}
\ce{Fe^{3+}} $3\text{d}^5$~: 5~; \ce{Mn^{2+}} $3\text{d}^5$~: 5~; \ce{Fe^{2+}}
$3\text{d}^6$~: 4~; \ce{Cu^{2+}} $3\text{d}^9$~: 1~; \ce{Zn^{2+}}
$3\text{d}^{10}$~: 0. Classement~: $\ce{Fe^{3+}} = \ce{Mn^{2+}} > \ce{Fe^{2+}}
> \ce{Cu^{2+}} > \ce{Zn^{2+}}$. Le fer(III) a un \omterm{def:b1:quantum-numbers:unpaired}{électron célibataire}
de plus que le fer(II), il est donc plus attiré~; le zinc(II) n'en a
aucun, il n'est donc pas attiré.
% ledger: gs:Fe2+, gs:Fe3+, gs:Mn2+, gs:Cu2+, gs:Zn2+
\end{solution}

\begin{solution}{pb:b1:quantum-numbers:1}
\textbf{1.} $E_1 = \qty{-13.6}{eV}$, $E_2 = \qty{-3.40}{eV}$,
$E_3 = \qty{-1.51}{eV}$, $E_4 = \qty{-0.85}{eV}$.
\textbf{2.} $\Delta E = 3.40 - 1.51 = \qty{1.89}{eV}$, $\lambda = 1240/1.89
= \qty{656}{nm}$~: rouge.
\textbf{3.} $4 \to 2$~: \qty{2.55}{eV}, \qty{486}{nm} (bleu-vert)~;
$5 \to 2$~: \qty{2.86}{eV}, \qty{434}{nm} (violet)~; $6 \to 2$~: \qty{3.022}{eV},
\qty{410}{nm} (violet).
\textbf{4.} $7 \to 2$~: \qty{3.12}{eV}, \qty{397}{nm}, déjà dans
l'ultraviolet~; les transitions plus hautes sont plus courtes encore, et
$3 \to 2$ est la plus grande longueur d'onde de Balmer. Les raies de
Lyman sont toutes ultraviolettes et celles de Paschen toutes
infrarouges~: seules quatre raies sont visibles.
\textbf{5.} $\infty \to 2$~: \qty{3.40}{eV}, $\lambda = \qty{365}{nm}$.
\textbf{6.} $\Delta E = 13.6 \times 3/4 = \qty{10.2}{eV}$, $\lambda =
\qty{122}{nm}$~: ultraviolet lointain, invisible à l'œil et absorbé par
le verre et l'air.
\textbf{7.} Le photon doit apporter au moins \qty{13.6}{eV}~: $\lambda \le
1240/13.6 = \qty{91.2}{nm}$.
\textbf{8.} $l = 0$ ($m_l = 0$), $l = 1$ ($m_l = -1, 0, 1$), $l = 2$
($m_l = -2, \dots, 2$)~: $1 + 3 + 5 = 9$ orbitales.
\textbf{9.} $\sum_{l=0}^{n-1}(2l+1) = n^2$ orbitales de deux électrons~:
$2n^2$~; pour $n = 4$, 32.
\textbf{10.} 1s, 2s, 2p, 3s, 3p, 4s, 3d, 4p, 5s, 4d, 5p, 6s, 4f, 5d, 6p,
7s, 5f, 6d, 7p.
\textbf{11.} Chaque période va de $n$s à $n$p~: 2, 8, 8, 18, 18, 32, 32.
\textbf{12.} 3d ($n + l = 5$) vient après 4s ($n + l = 4$), qui ouvre
la période 4.
\textbf{13.} $Z = 2, 10, 18, 36, 54, 86, 118$.
\textbf{14.} $[\ce{Ar}]\,3\text{d}^6 4\text{s}^2$~; 3d~: \omorbs{ud,u,u,u,u},
4 \omterm{def:b1:quantum-numbers:unpaired}{électrons célibataires}.
\textbf{15.} \ce{Fe^{2+}} $[\ce{Ar}]\,3\text{d}^6$~: 4 célibataires~;
\ce{Fe^{3+}} $[\ce{Ar}]\,3\text{d}^5$~: 5 célibataires, une \omterm{def:b1:quantum-numbers:orbital}{sous-couche} à
moitié remplie.
\textbf{16.} Prévu $3\text{d}^4 4\text{s}^2$~: 4 célibataires~; observé
$3\text{d}^5 4\text{s}^1$~: 6 célibataires (5 en 3d, 1 en 4s).
\textbf{17.} \ce{Cu} $[\ce{Ar}]\,3\text{d}^{10} 4\text{s}^1$, \ce{Cu+}
$[\ce{Ar}]\,3\text{d}^{10}$, \ce{Cu^{2+}} $[\ce{Ar}]\,3\text{d}^9$.
% ledger: gs:Cu, gs:Cu+, gs:Cu2+
\textbf{18.} \ce{Cu^{2+}}, avec un \omterm{def:b1:quantum-numbers:unpaired}{électron célibataire}~; \ce{Cu+} n'en a
aucun.
\textbf{19.} \ce{Sc^{3+}} $[\ce{Ar}]$~: 0~; \ce{Ti^{3+}} $[\ce{Ar}]\,3\text{d}^1$~:
1~; \ce{Mn^{2+}} $[\ce{Ar}]\,3\text{d}^5$~: 5. C'est \ce{Mn^{2+}} qui en a
le plus.
\textbf{20.} Trois électrons par orbitale~; $3(2l + 1)$ par \omterm{def:b1:quantum-numbers:orbital}{sous-couche}~;
$3n^2$ par couche.
\textbf{21.} 1s~: 3~; 2s~: 3~; 2p~: 9~; 3s~: 3~; 3p~: 9.
\textbf{22.} La première période se ferme quand 1s est pleine~: $Z = 3$.
\textbf{23.} $Z = 4$, un électron en 2s au-delà du cœur de gaz noble.
\textbf{24.} 2s et 2p~: $3 + 9 = 12$ éléments.
\textbf{25.} $3 + 12 = 15$~: le deuxième gaz noble du monde à trois
spins a \textbf{$Z = 15$}.
\end{solution}
